\clearpage\fancyhead[L]{\small COMET.}\SourceLeaf{387}
\begin{center}\textbf{COMET.}\par\vspace{4pt}\hrule\vspace{4pt}\fontsize{10}{12}\selectfont Göttingische gelehrte Anzeigen. No. 45. p. 441. 20 March 1815.\par\vspace{4pt}\hrule\vspace{4pt}\end{center}

\par\indent In a letter of 7 March to Professor \textsc{Gauss}, Dr. \textsc{Olbers} reported that he had discovered a new comet in Perseus on 6 March. On 6 March at $10^{\mathrm u}\,55^{\mathrm m}$ mean time its right ascension was $49^\circ\,7'$ and its northern declination $32^\circ\,7'$. On 7 March at $7^{\mathrm u}\,40^{\mathrm m}$ its right ascension was $49^\circ\,22'$ and its northern declination $32^\circ\,32'$. On this basis, the comet can be expected to remain visible for quite some time. It is very small, has a poorly defined nucleus and a very faint, transparent nebulosity, and is similar in brightness and shape to the second comet of 1811 as it appeared in February 1812.

\par\indent The search for the new comet at the observatory here began as soon as the letter arrived on 13 March. For only a few minutes the sky cleared enough for the comet to be seen between 159 and 164 in Perseus (according to \textsc{Bode}’s catalogue). But the rain clouds returned so quickly that even an estimate was impossible.
